\section{Operations over edges in force}
\label{sec:queries}

Neighbors, path, communities and answer read the graph through one function: the
edges each pair holds at a point, resolved by the rule of \S\ref{sec:model}.
A walk, a path and a partition therefore cannot disagree about what the graph
was at a point, and each of them, asked at a past $\tk$, repeats the answer it
would have given then.

\subsection{Walks and their bounds}

\shipped{} \op{neighbors} walks from up to 256 seeds, out along edges, in against
them, or both, for up to 32 hops, and returns at most 10{,}000 entities ordered by
hop count and then by key. A hop reads each entity's pairs that have held an
edge, including their properties, since a retraction is a property that can end
an edge, and keeps what each pair holds at the point. An edge that has ended, or
has been superseded or retracted, is not crossed, and a property is never
crossed. Each organization's file is served through one connection, so a long
walk would block every other write for that organization; the bounds limit the
work as well as the answer, and a response that reached one of them says so.

\subsection{Shortest path and confidence decay}

\shipped{} \op{path} is the same walk stopped at the level that first reaches a
goal and read back along the hop that first reached each node. Levels are read
in order and candidate hops are sorted by source, relation and target before
they are claimed, so the path is a shortest path in hops and ties are broken by
key rather than by the engine's row order. It may be restricted to a list of
relations.

The response carries the product of the edges' confidences and the edge of
lowest confidence. Both are omitted when any edge on the path states no
confidence, since an absent confidence reads as zero and a product over values
nobody supplied measures nothing. The product is a summary of the path found. It
is not the most confident path: selection is by hop count, and a longer path of
higher confidence is not considered. The product also treats the edges'
confidences as independent, which the store has no way to check.

\subsection{Communities}

\shipped{} \op{communities} partitions the edges in force at a point, with
direction dropped and self-loops excluded. Two entities are joined with weight
equal to the number of distinct (relation, direction) edges between them, so an
edge asserted by three sources counts once. The algorithm is Louvain
modularity optimisation~\cite{blondel2008louvain,newman2004modularity}, made
deterministic: nodes are visited in key order; a node moves to the neighbouring
community of greatest modularity gain, stays where it is on a tie with its own
community, and otherwise goes to the lower-numbered one; communities are
numbered by size and then by lowest key. Each pass's partition is reported as a
level, coarsest first, with its modularity.

Louvain can return a community whose members are not connected to each other;
Traag et al.\ describe the defect and give the Leiden refinement as a
remedy~\cite{traag2019leiden}. The store takes a smaller step: after each pass,
every community is split into its connected components in the original graph.

\begin{proposition}[Splitting by components does not lower modularity]
Let community $c$ be the union of two sets $c_1, c_2$ with no edge between them.
Replacing $c$ by $c_1$ and $c_2$ changes modularity by
$2\,d_{c_1} d_{c_2} / (2m)^2 \ge 0$, where $d$ is total degree and $m$ total
edge weight.
\end{proposition}
\begin{proof}
Modularity is $Q = \sum_c \big[ w_c / 2m - (d_c / 2m)^2 \big]$ with $w_c$ twice
the weight inside $c$. With no edge between $c_1$ and $c_2$,
$w_c = w_{c_1} + w_{c_2}$ and $d_c = d_{c_1} + d_{c_2}$, so the first term is
unchanged and the second falls by $2 d_{c_1} d_{c_2}/(2m)^2$.
\end{proof}

A component of a finer level lies inside one community of the coarser level and
is connected, so it lies inside one component of that community: the split keeps
the levels nested. The split guarantees connected communities; it does not
provide Leiden's further guarantee that communities are well connected, and
Leiden itself is not implemented. A partition reads at most 100{,}000 edge rows,
and a graph over that bound is refused rather than partitioned from part of
itself. Nothing is stored: a partition is recomputed on every call from the rows.

\subsection{An answer over the whole graph}

\shipped{} \op{answer} addresses a question whose answer's location is unknown,
so no walk can start. It follows the map-reduce shape of
GraphRAG~\cite{edge2024graphrag} over the partition at the requested point, with
two differences: nothing is summarized or stored in advance, and every citation
is checked.

\begin{enumerate}[leftmargin=*,itemsep=2pt]
\item The coarsest level's communities are ordered: those containing an entity
  that a full-text match of the question's words finds come first, ranked by
  BM25~\cite{robertson2009bm25}, and the rest follow by size.
\item Communities are read in that order into batches of in-force facts, each
  fact labelled $[i]$: at most 12\,KiB of facts per batch, eight batches, 2{,}048
  members in all, 256 per community and 64 assertions per member, with members
  the question names read first and then the most connected.
\item Each batch is one model call that returns at most ten findings, each citing
  the labels it rests on. A label outside the batch is dropped, and a finding
  that cites nothing it was given, or scores itself zero, is discarded.
\item One further call receives the findings, ordered by score, and returns an
  answer citing the findings it rests on.
\end{enumerate}

The response lists the assertion identifiers and community numbers behind each
cited finding. Each identifier is a row that was given to the model at the
point, so the model can omit a citation but cannot invent one. An empty answer
is returned as ungrounded. A non-empty answer that cites no finding is refused
with an error rather than returned. A failure of any model call fails the whole
answer, since an answer reduced from the batches that succeeded would read as
complete. The response reports how many communities were read and skipped,
whether any bound clipped the facts, and how many model calls were made; the
calls are billed to the calling organization through the platform's model
client, which checks the balance before each call.

These checks establish provenance of the citations, not entailment. Nothing
verifies that a cited fact supports the sentence that cites it, and the answer
depends on the model the platform routes to.
